\documentclass[11pt]{article}
\usepackage{amsmath}
\usepackage[margin=1in]{geometry}
\usepackage{booktabs}
\usepackage{longtable}
\usepackage{textcomp}
\usepackage{caption}
\usepackage[backend=biber,style=authoryear,natbib=true]{biblatex}
\addbibresource{references.bib}

\title{Supplementary Table~S1: DAYCENT \texttt{crop.100} parameters for \texttt{MISC} and \texttt{SG3} entries}
\author{}
\date{}

\begin{document}
\maketitle

\begin{center}
\textbf{Table~S1.} DAYCENT \texttt{crop.100} parameters for the \texttt{MISC} (miscanthus proxy) and
\texttt{SG3} (switchgrass) crop entries used in this study. Both entries were derived from the
warm-season grass (G2) base parameterization \citep{WangDou2017DAYCENT,WangDou2020AgronomyJ}.
The only differing parameter is PRDX(1) (highlighted). All parameters were verified by direct
inspection of the \texttt{crop.100} file (Yong Wang, archived simulation, DDcentEVI revision 279).
\end{center}

\vspace{1em}

\begin{longtable}{llllp{6cm}}
\toprule
\textbf{Parameter} & \textbf{MISC} & \textbf{SG3} & \textbf{G2 base} & \textbf{Description} \\
\midrule
\endfirsthead
\multicolumn{5}{c}{\tablename\ \thetable{} -- continued} \\
\toprule
\textbf{Parameter} & \textbf{MISC} & \textbf{SG3} & \textbf{G2 base} & \textbf{Description} \\
\midrule
\endhead
\bottomrule
\endlastfoot
\textbf{PRDX(1)} & \textbf{3.5} & \textbf{2.75} & 1.5 & \textbf{Potential production slope, above-ground (g\,C\,m$^{-2}$); only differing parameter} \\
PRDX(2) & 3.0 & 3.0 & 1.5 & Potential production slope, below-ground \\
PPDF(1) & 30.0 & 30.0 & 30.0 & Optimum temperature for plant production (°C) \\
PPDF(2) & 45.0 & 45.0 & 45.0 & Maximum temperature for plant production (°C) \\
PPDF(3) & 1.0 & 1.0 & 1.0 & Growth response minimum (fraction) \\
PPDF(4) & 2.5 & 2.5 & 3.0 & Growth response shape \\
BIOFLG & 1 & 1 & 1 & Biomass allocation flag \\
BIOK5 & 60.0 & 60.0 & 60.0 & Minimum temperature, above-ground biomass \\
BIOMAX & 200.0 & 200.0 & 200.0 & Maximum standing dead biomass (g\,C\,m$^{-2}$) \\
PLTMRF & 1.0 & 1.0 & 1.0 & Plant maturity factor \\
FULCAN & 100.0 & 100.0 & 100.0 & Full canopy factor \\
FRTCINDX & 1 & 1 & 1 & Root C allocation index \\
FRTC(1) & 0.75 & 0.75 & 0.75 & Root C allocation fraction 1 \\
FRTC(2) & 0.40 & 0.40 & 0.35 & Root C allocation fraction 2 \\
FRTC(3) & 90.0 & 90.0 & 90.0 & Root growth days \\
FRTC(4) & 0.20 & 0.20 & 0.20 & Root C allocation fraction 4 \\
FRTC(5) & 0.10 & 0.10 & 0.10 & Root C allocation fraction 5 \\
CFRTCN(1) & 0.50 & 0.50 & 0.60 & Fine root C fraction, nitrogen-limited \\
CFRTCN(2) & 0.30 & 0.30 & 0.45 & Fine root C fraction, nitrogen-limited 2 \\
CFRTCW(1) & 0.60 & 0.60 & 0.60 & Fine root C fraction, water-limited \\
CFRTCW(2) & 0.30 & 0.30 & 0.30 & Fine root C fraction, water-limited 2 \\
HIMAX & 0.02 & 0.02 & 0.02 & Maximum harvest index \\
HIWSF & 0.50 & 0.50 & 0.00 & Harvest water stress factor \\
FALLRT & 0.01 & 0.01 & 0.15 & Fall senescence rate \\
RDRM & 0.20 & 0.20 & 0.20 & Root death rate, mature \\
RDRJ & 0.40 & 0.40 & 0.40 & Root death rate, juvenile \\
\bottomrule
\end{longtable}

\textit{Note:} All 112 parameters in the \texttt{crop.100} entries for \texttt{MISC} and \texttt{SG3} were compared by direct inspection of the archived parameter file (Yong Wang, DDcentEVI revision 279). Only PRDX(1) differs between the two entries. The table above lists all parameters that deviate from the G2 warm-season grass base parameterization, plus the primary biological controls on biomass production and root allocation; all remaining parameters (not shown) are set to identical DAYCENT system defaults for both \texttt{MISC} and \texttt{SG3}. The G2 column shows the warm-season grass base parameterization from which both crop entries were derived; parameters differing from the G2 base were updated identically for both \texttt{MISC} and \texttt{SG3} in the Wang bioenergy crop calibration.

\printbibliography

\end{document}
